\documentclass[10pt,letterpaper]{article}
\usepackage[utf8]{inputenc}
\usepackage[T1]{fontenc}
\usepackage{mathptmx}
\renewcommand{\ttdefault}{txtt}
\usepackage[left=1.9cm,right=1.9cm,top=3.2cm,bottom=2.5cm,headheight=1.7cm,headsep=0.4cm,footskip=0.9cm]{geometry}
\usepackage{xcolor}
\usepackage{graphicx}
\usepackage[font={footnotesize,it},labelsep=colon,hypcap=false]{caption}
\usepackage{fancyhdr}
\usepackage{eso-pic}
\usepackage{titlesec}
\usepackage{listings}
\usepackage[hidelinks]{hyperref}

\linespread{1.1}
\setlength{\parindent}{0pt}
\setlength{\parskip}{6pt}
\definecolor{gris}{gray}{0.35}
\setlength{\fboxsep}{0pt}
\setlength{\fboxrule}{0.4pt}
\captionsetup[figure]{name=Imagen}
\titleformat*{\section}{\normalsize\bfseries}
\titlespacing*{\section}{0pt}{14pt}{6pt}

\pagestyle{fancy}
\fancyhf{}
\renewcommand{\headrulewidth}{0pt}
\fancyhead[L]{\includegraphics[height=1.4cm]{logo_usm}}
\fancyhead[R]{\footnotesize UNIVERSIDAD TÉCNICA FEDERICO SANTA MARÍA\\DEPARTAMENTO DE INDUSTRIAS\\CAMPUS SANTIAGO VITACURA\\SEGUNDO SEMESTRE 2026}
\fancyfoot[R]{\footnotesize\thepage}

\AddToShipoutPictureBG{\AtPageLowerLeft{\setlength{\unitlength}{1bp}\color[gray]{0.6}\put(24,24){\framebox(564,744){}}}}

\definecolor{clave}{HTML}{1A8BCB}
\definecolor{funcion}{HTML}{A8A898}
\definecolor{texto}{HTML}{E08A14}
\definecolor{alias}{HTML}{A0422E}
\definecolor{codigo}{HTML}{212121}

\makeatletter
\newcommand{\numero}[1]{\ifnum\lst@mode=\lst@Pmode\textcolor{texto}{#1}\else#1\fi}
\makeatother

\lstset{
  morestring=[b]',
  moredelim=[s][\color{alias}]{`}{`},
  keywords={SELECT,FROM,WHERE,AND,OR,NOT,IS,NULL,IN,BETWEEN,ORDER,BY,GROUP,HAVING,LIMIT,DESC,ASC,AS,LEFT,JOIN,ON,DISTINCT,SET,UPDATE,DELETE,INSERT,INTO,VALUES,AVG},
  keywords=[2]{COUNT,SUM,MIN,MAX,COALESCE},
  basicstyle=\ttfamily\footnotesize\color{codigo},
  keywordstyle=\bfseries\color{clave},
  keywordstyle=[2]\color{funcion},
  stringstyle=\color{texto},
  frame=tb,
  columns=fullflexible,
  keepspaces=true,
  upquote=true,
  showstringspaces=false,
  literate={á}{{\'a}}1 {é}{{\'e}}1 {í}{{\'i}}1 {ó}{{\'o}}1 {ú}{{\'u}}1 {ñ}{{\~n}}1
           {0}{{\numero{0}}}1 {1}{{\numero{1}}}1 {2}{{\numero{2}}}1 {3}{{\numero{3}}}1 {4}{{\numero{4}}}1 {5}{{\numero{5}}}1 {6}{{\numero{6}}}1 {7}{{\numero{7}}}1 {8}{{\numero{8}}}1 {9}{{\numero{9}}}1
}

\begin{document}

\begin{center}
\vspace*{4cm}
{\fontsize{26}{32}\selectfont\bfseries\textsc{Laboratorio 5}\par}
\vspace{10pt}
{\color{lightgray}\rule{0.75\linewidth}{0.5pt}}\par
\vspace{2pt}
{\itshape Entregable SQL (2/2): base de datos Andes Bike\par}
{\color{lightgray}\rule{0.75\linewidth}{0.5pt}}\par
\vspace{12pt}
{\fontsize{22}{28}\selectfont\bfseries\textsc{Sistemas de información para la gestión.}\par}
\vspace{4.5cm}
{\color{gris}\itshape Daniela Román Schuffenegger, rol 202160649-1, \href{mailto:daniela.roman@usm.cl}{\underline{daniela.roman@usm.cl}}\\
Paralelo 101\par}
\vspace{1.8cm}
{\color{gris}\itshape Profesor: José Luis Sáez.\\
Ayudantes: Valeria Bobadilla, Catalina Bodaleo, Luis Contreras, Benjamín Jara.\\
Fecha: 10 de noviembre de 2026\par}
\end{center}
\vspace{0.8cm}
\textbf{Repositorio Github: \href{https://github.com/DanielaRoman7/ICN292-Lab5-Roman-Daniela}{https://github.com/DanielaRoman7/ICN292-Lab5-Roman-Daniela}}
\newpage

\section*{Consulta 1: ventas de cada vendedor}

\begin{minipage}{\linewidth}
\begin{lstlisting}
SELECT v.nombre AS `Vendedor`, COUNT(ve.id_venta) AS `Cantidad de ventas`
FROM vendedores v
LEFT JOIN ventas ve ON v.id_vendedor = ve.id_vendedor
GROUP BY v.id_vendedor, v.nombre
ORDER BY `Cantidad de ventas` DESC, v.nombre;
\end{lstlisting}
\centering
\fcolorbox{gray}{white}{\includegraphics[width=0.31\linewidth]{capturas/consulta01}}
\captionof{figure}{Cantidad de ventas por vendedor}
\end{minipage}

\section*{Consulta 2: ventas y unidades por tienda}

\begin{minipage}{\linewidth}
\begin{lstlisting}
SELECT t.nombre AS `Tienda`,
       COUNT(ve.id_venta) AS `Ventas`,
       SUM(ve.unidades) AS `Unidades vendidas`
FROM tiendas t
LEFT JOIN vendedores v ON t.id_tienda = v.id_tienda
LEFT JOIN ventas ve ON v.id_vendedor = ve.id_vendedor
GROUP BY t.id_tienda, t.nombre
ORDER BY `Unidades vendidas` DESC;
\end{lstlisting}
\centering
\fcolorbox{gray}{white}{\includegraphics[width=0.39\linewidth]{capturas/consulta02}}
\captionof{figure}{Ventas y unidades vendidas por tienda}
\end{minipage}

\section*{Consulta 3: tiendas sin vendedores}

\begin{minipage}{\linewidth}
\begin{lstlisting}
SELECT t.nombre AS `Tienda`, t.comuna AS `Comuna`
FROM tiendas t
LEFT JOIN vendedores v ON t.id_tienda = v.id_tienda
WHERE v.id_vendedor IS NULL;
\end{lstlisting}
\centering
\fcolorbox{gray}{white}{\includegraphics[width=0.24\linewidth]{capturas/consulta03}}
\captionof{figure}{Tiendas sin vendedores asignados}
\end{minipage}

\section*{Consulta 4: modelos sin ventas en junio de 2025}

\begin{minipage}{\linewidth}
\begin{lstlisting}
SELECT b.modelo AS `Modelo`, b.categoria AS `Categoría`
FROM bicicletas b
LEFT JOIN ventas ve ON b.id_bicicleta = ve.id_bicicleta
    AND ve.fecha BETWEEN '2025-06-01' AND '2025-06-30'
WHERE ve.id_venta IS NULL
ORDER BY b.modelo;
\end{lstlisting}
\centering
\fcolorbox{gray}{white}{\includegraphics[width=0.21\linewidth]{capturas/consulta04}}
\captionof{figure}{Modelos sin ventas en junio de 2025}
\end{minipage}

\bigskip

Puse la condición del mes en el \lstinline|ON| del \lstinline|LEFT JOIN|, así solo se unen las ventas de junio y las bicicletas que no tuvieron ninguna quedan con \lstinline|NULL|. Si la hubiera puesto en el \lstinline|WHERE|, ese filtro habría eliminado justamente las filas con \lstinline|NULL| y la consulta no habría devuelto ningún modelo.

\section*{Consulta 5: vendedores sin ventas}

\begin{minipage}{\linewidth}
\begin{lstlisting}
SELECT v.nombre AS `Vendedor`, t.nombre AS `Tienda`
FROM vendedores v
LEFT JOIN tiendas t ON v.id_tienda = t.id_tienda
LEFT JOIN ventas ve ON v.id_vendedor = ve.id_vendedor
WHERE ve.id_venta IS NULL
ORDER BY v.nombre;
\end{lstlisting}
\centering
\fcolorbox{gray}{white}{\includegraphics[width=0.29\linewidth]{capturas/consulta05}}
\captionof{figure}{Vendedores sin ventas y su tienda}
\end{minipage}

\section*{Consulta 6: vendedores de una misma tienda}

\begin{minipage}{\linewidth}
\begin{lstlisting}
SELECT a.nombre AS `Vendedor 1`,
       a.fecha_ingreso AS `Ingreso 1`,
       b.nombre AS `Vendedor 2`,
       b.fecha_ingreso AS `Ingreso 2`,
       t.nombre AS `Tienda`
FROM vendedores a
JOIN vendedores b ON a.id_tienda = b.id_tienda
    AND a.fecha_ingreso < b.fecha_ingreso
JOIN tiendas t ON a.id_tienda = t.id_tienda
ORDER BY t.nombre, a.fecha_ingreso, b.fecha_ingreso;
\end{lstlisting}
\centering
\fcolorbox{gray}{white}{\includegraphics[width=0.66\linewidth]{capturas/consulta06}}
\captionof{figure}{Pares de vendedores de la misma tienda}
\end{minipage}

\section*{Consulta 7: tiendas sobre \$2.000.000 en Montaña y Ruta}

\begin{minipage}{\linewidth}
\begin{lstlisting}
SELECT t.nombre AS `Tienda`, SUM(ve.unidades * b.precio) AS `Monto Montaña y Ruta`
FROM tiendas t
JOIN vendedores v ON t.id_tienda = v.id_tienda
JOIN ventas ve ON v.id_vendedor = ve.id_vendedor
JOIN bicicletas b ON ve.id_bicicleta = b.id_bicicleta
WHERE b.categoria IN ('Montaña', 'Ruta')
GROUP BY t.id_tienda, t.nombre
HAVING SUM(ve.unidades * b.precio) > 2000000
ORDER BY `Monto Montaña y Ruta` DESC;
\end{lstlisting}
\centering
\fcolorbox{gray}{white}{\includegraphics[width=0.33\linewidth]{capturas/consulta07}}
\captionof{figure}{Tiendas sobre \$2.000.000 en Montaña y Ruta}
\end{minipage}

\section*{Consulta 8: bicicletas sobre el precio promedio}

\begin{minipage}{\linewidth}
\textbf{Primera sentencia: precio promedio del catálogo}
\begin{lstlisting}
SELECT AVG(precio) AS `Precio promedio`
FROM bicicletas;
\end{lstlisting}
\centering
\fcolorbox{gray}{white}{\includegraphics[width=0.15\linewidth]{capturas/consulta08a}}
\captionof{figure}{Precio promedio del catálogo}
\end{minipage}

\bigskip

\begin{minipage}{\linewidth}
\textbf{Segunda sentencia: bicicletas que superan ese promedio}
\begin{lstlisting}
SELECT modelo AS `Modelo`, precio AS `Precio`
FROM bicicletas
WHERE precio > (SELECT AVG(precio) FROM bicicletas)
ORDER BY precio DESC;
\end{lstlisting}
\centering
\fcolorbox{gray}{white}{\includegraphics[width=0.23\linewidth]{capturas/consulta08b}}
\captionof{figure}{Bicicletas sobre el precio promedio}
\end{minipage}

\bigskip

El precio promedio del catálogo es \$657.490 y lo obtuve con la subconsulta \lstinline|SELECT AVG(precio) FROM bicicletas|, que promedia el precio de las 8 bicicletas. En la primera sentencia lo muestro por separado y en la segunda lo uso dentro del \lstinline|WHERE| para filtrar.

\section*{Consulta 9: vendedores sobre el promedio de unidades}

\begin{minipage}{\linewidth}
\begin{lstlisting}
SELECT v.nombre AS `Vendedor`, AVG(ve.unidades) AS `Promedio de unidades`
FROM vendedores v
JOIN ventas ve ON v.id_vendedor = ve.id_vendedor
GROUP BY v.id_vendedor, v.nombre
HAVING AVG(ve.unidades) > (SELECT AVG(unidades) FROM ventas)
ORDER BY `Promedio de unidades` DESC, v.nombre;
\end{lstlisting}
\centering
\fcolorbox{gray}{white}{\includegraphics[width=0.31\linewidth]{capturas/consulta09}}
\captionof{figure}{Vendedores sobre el promedio de unidades de la cadena}
\end{minipage}

\section*{Consulta 10: monto total por vendedor}

\begin{minipage}{\linewidth}
\begin{lstlisting}
SELECT v.nombre AS `Vendedor`, SUM(ve.unidades * b.precio) AS `Monto total`
FROM vendedores v
LEFT JOIN ventas ve ON v.id_vendedor = ve.id_vendedor
LEFT JOIN bicicletas b ON ve.id_bicicleta = b.id_bicicleta
GROUP BY v.id_vendedor, v.nombre
ORDER BY `Monto total` DESC;
\end{lstlisting}
\centering
\fcolorbox{gray}{white}{\includegraphics[width=0.25\linewidth]{capturas/consulta10}}
\captionof{figure}{Monto total vendido por vendedor}
\end{minipage}

\bigskip

Fernanda Lillo y Rubén Castro aparecen con \lstinline|NULL| y no con 0 porque no tienen ventas: el \lstinline|LEFT JOIN| los mantiene, pero les deja en \lstinline|NULL| las columnas que vienen de ventas y bicicletas, y \lstinline|SUM| sobre puros \lstinline|NULL| devuelve \lstinline|NULL|. Para que mostraran 0 habría que envolver la suma con \lstinline|COALESCE(SUM(ve.unidades * b.precio), 0)|.

\end{document}
